%
% Documento: Resumo (Inglês)
%


\begin{center}
\imprimirtitleenglish
\end{center}

\begin{resumo}[Abstract]
            This research aims to contribute to the determination of a fundamental parameter in astronomy: the mass of young stars in star clusters, with the potential to enrich the understanding of stellar rotation. The aim is to estimate stellar masses in clusters such as the Orion Nebula (ONC), Pleiades, h Persei, Upper Scorpius and NGC 2516 using a combination of traditional methods, such as fitting isochrones and mass-luminosity relations, and machine learning techniques. Theoretical isochrones and evolutionary tracks were applied to estimate masses and provide information on cluster ages and stellar evolution. This study takes a descriptive approach to investigating stellar mass calculation methods, using statistical tools to assess the quality of the estimates. The sample is non-probabilistic, with a focus on young low-mass stars, and data was collected from several catalogs, including GAIA and TESS. Absolute magnitudes were calculated from photometry, and individual interstellar extinction was considered for the youngest clusters. Stellar mass was estimated by means of isochrones and evolutionary models, using machine learning techniques, including K-fold cross-validation and grid search optimization, with various regression techniques such as Random Forests, Support Vector Regression, and Bayesian Regression. The results were evaluated with a set of statistical metrics, such as RMSE, AIC and Coefficient of Determination, compared with mass values recently obtained in the literature, demonstrating the effectiveness of the mass-luminosity relationship, which had a good agreement with most clusters, especially those with a well-defined main sequence, such as the Pleiades, and of the isochrones adjustment for clusters with more objects in the main pre-sequence, such as the ONC, although with inherent uncertainties and limitations. The integration of machine learning techniques has shown promise for increasing the accuracy and reliability of stellar mass estimates, pointing to future directions for astrophysical research.

            \vspace*{12 pt}

            
            \textbf{Keywords}: Stellar Parameters: Mass. Computational Modeling. Statistics.

\end{resumo}

	
